%% P24 --- Sums of random variables: the central limit theorem, concentration and Monte Carlo
%%
%% Stub. The brief below is the contract for this program: what it must argue,
%% and what it must buy the reader. Writing the program means deleting the
%% \programstub{} block and replacing it with the program. If the brief turns
%% out to be wrong once the mathematics has been worked, change it and record
%% the contradiction in CLAUDE.md under *Resolved questions*.

\program{Sums of random variables: the central limit theorem, concentration and Monte Carlo}
\label{prog:P24}

\programstub{%
Argument: variances of independent quantities add; averages concentrate; the
error of an average falls as \code{1/sqrt(n)} and that single rate governs
an astonishing amount of practice. Payoff: the derivation of the
\code{1/sqrt(d\_k)} in attention. A dot product of two
\code{d\_k}-dimensional vectors with independent unit-variance entries has
variance \code{d\_k}, so its standard deviation grows as \code{sqrt(d\_k)};
feed that into a softmax and it saturates, and the gradient dies; divide by
\code{sqrt(d\_k)} and the logits have unit variance regardless of head size.
The scaling is a variance correction and nothing else, and the reader
derives it here and measures it in P31. Also: Monte Carlo error is the same
\code{1/sqrt(n)}, so the number of samples needed to halve an error bar is
four times as many --- the fact that prices every evaluation run in the
book.

\medskip
\textbf{Planned length:} 55 frames.
}
